\documentclass[12pt]{article}
\usepackage{amsmath}
\usepackage{amssymb}
\usepackage{geometry}
\usepackage{xcolor}
\usepackage{tcolorbox}
\geometry{a4paper, margin=1in}

\title{\textbf{Homework 4.5: The Superposition Principle and Undetermined Coefficients Revisited}}
\author{Tashfeen Omran}
\date{September 2026}

\begin{document}

\maketitle
\tableofcontents
\newpage

\section{Part 1: Comprehensive Introduction, Context, and Prerequisites}

\subsection{Core Concepts}
In this chapter, we expand our toolkit for solving nonhomogeneous linear second-order differential equations of the form:
\[ a y'' + b y' + c y = f(t) \]
where $f(t)$ is a non-zero forcing function. The general solution to this equation is always the sum of two distinct parts:
\[ y(t) = y_h(t) + y_p(t) \]
Here, $y_h(t)$ is the general solution to the associated homogeneous equation ($a y'' + b y' + c y = 0$), and $y_p(t)$ is \textit{any} particular solution to the nonhomogeneous equation. 

The two major concepts introduced in this section are:
\begin{itemize}
    \item \textbf{The Superposition Principle:} If a forcing function $f(t)$ is made up of a sum of simpler functions, $f(t) = f_1(t) + f_2(t)$, we can find particular solutions $y_{p1}$ and $y_{p2}$ for each piece individually. The total particular solution is simply $y_p = y_{p1} + y_{p2}$.
    \item \textbf{Undetermined Coefficients (Revisited):} A systematic method for guessing the form of $y_p(t)$ when $f(t)$ consists of polynomials, exponentials, sines, cosines, or finite products/sums of these. We guess a form with unknown constants (coefficients) and plug it into the differential equation to solve for them.
\end{itemize}

\subsection{Intuition and Derivation}
The Superposition Principle is a direct consequence of the \textit{linearity} of the differential operator. Let $L[y] = a y'' + b y' + c y$. Because derivatives are linear, $L[y_1 + y_2] = L[y_1] + L[y_2]$. Therefore, if $L[y_1] = f_1(t)$ and $L[y_2] = f_2(t)$, then $L[y_1 + y_2] = f_1(t) + f_2(t)$. This allows us to divide and conquer complex forcing functions.

The Method of Undetermined Coefficients works because the derivatives of polynomials, exponentials, sines, and cosines result in functions of the same type. If $f(t) = e^{2t}$, it is highly intuitive that $y_p(t)$ must also involve $e^{2t}$, because only an exponential can differentiate into an exponential and balance the equation.

\subsection{Historical Context and Motivation}
The mathematical foundations for linear differential equations were laid in the 18th century by Leonhard Euler and Joseph-Louis Lagrange. They were motivated by classical mechanics, specifically the study of vibrations and planetary orbits. In real-world physical systems, an object (like a mass on a spring or a bridge) is often subjected to multiple external forces simultaneously (e.g., wind, traffic, and seismic activity). The Superposition Principle was a monumental breakthrough because it proved that the total response of a linear system is simply the sum of the responses to each individual force, making complex physical modeling mathematically tractable.

\subsection{Key Formulas}
\begin{itemize}
    \item \textbf{General Solution:} $y(t) = c_1 y_1(t) + c_2 y_2(t) + y_p(t)$
    \item \textbf{Form of the Particular Solution:} For a forcing term $f(t) = P_m(t) e^{\alpha t} \cos(\beta t) + Q_n(t) e^{\alpha t} \sin(\beta t)$, the guess is:
    \[ y_p(t) = t^s \left[ (A_k t^k + \dots + A_0) e^{\alpha t} \cos(\beta t) + (B_k t^k + \dots + B_0) e^{\alpha t} \sin(\beta t) \right] \]
    where $k = \max(m, n)$, and $s$ is the number of times the complex number $\alpha + i\beta$ is a root of the characteristic equation $a r^2 + b r + c = 0$. (If it is not a root, $s=0$).
\end{itemize}

\subsection{Prerequisites}
To succeed in this section, you must be proficient in:
\begin{enumerate}
    \item Solving homogeneous linear equations by finding the roots of the characteristic polynomial $a r^2 + b r + c = 0$.
    \item The Product Rule and Chain Rule for differentiation.
    \item Solving systems of linear algebraic equations (e.g., 2 equations with 2 unknowns).
\end{enumerate}

\newpage
\section{Part 2: Detailed Homework Solutions}
\textit{Note: To maximize learning efficiency, a strategic selection of problems has been chosen. This selection covers every distinct problem archetype, algebraic trick, and conceptual trap present in the chapter without redundant repetition.}

\subsection{Problem 1(c): Applying the Superposition Principle}
\textbf{Problem:} Given that $y_1(t) = \cos t$ is a solution to $y'' - y' + y = \sin t$ and $y_2(t) = e^{2t}/3$ is a solution to $y'' - y' + y = e^{2t}$, use the superposition principle to find a solution to $y'' - y' + y = 4\sin t + 18e^{2t}$.

\textbf{Solution:}
Let the differential operator be $L[y] = y'' - y' + y$.
We are given:
\[ L[y_1] = \sin t \quad \text{and} \quad L[y_2] = e^{2t} \]
We want to find a function $y_p$ such that:
\[ L[y_p] = 4\sin t + 18e^{2t} \]
By the linearity of the operator (Superposition Principle), we know that $L[k_1 y_1 + k_2 y_2] = k_1 L[y_1] + k_2 L[y_2]$.
Therefore, we can construct our solution by scaling the given solutions:
\[ y_p(t) = 4y_1(t) + 18y_2(t) \]
Substitute the given functions for $y_1$ and $y_2$:
\[ y_p(t) = 4(\cos t) + 18\left(\frac{e^{2t}}{3}\right) \]
Simplify the expression:
\[ y_p(t) = 4\cos t + 6e^{2t} \]
\textbf{Final Answer:} $y_p(t) = 4\cos t + 6e^{2t}$

\subsection{Problem 11: Applicability of Undetermined Coefficients}
\textbf{Problem:} Determine whether the method of undetermined coefficients can be applied to find a particular solution of the equation: $y'' - 6y' - 4y = 4\sin 3t - e^{3t}t^2 + 1/t$. Do not solve.

\textbf{Solution:}
The method of undetermined coefficients requires the nonhomogeneous forcing function $f(t)$ to be a linear combination of functions of the form $P(t)e^{\alpha t}\cos(\beta t)$ or $P(t)e^{\alpha t}\sin(\beta t)$, where $P(t)$ is a polynomial with non-negative integer exponents.
Let's analyze the terms of our forcing function $f(t) = 4\sin 3t - e^{3t}t^2 + t^{-1}$:
\begin{itemize}
    \item $4\sin 3t$: This is valid ($\alpha=0, \beta=3, P(t)=4$).
    \item $-e^{3t}t^2$: This is valid ($\alpha=3, \beta=0, P(t)=-t^2$).
    \item $1/t = t^{-1}$: This is \textbf{invalid}. The exponent is negative, meaning it is not a polynomial. The derivatives of $t^{-1}$ generate an infinite sequence of new linearly independent terms ($t^{-2}, t^{-3}, \dots$), making it impossible to form a closed finite guess for $y_p$.
\end{itemize}
\textbf{Final Answer:} No, the method cannot be applied because of the $1/t$ term.

\subsection{Problem 19: Finding a General Solution from Scratch}
\textbf{Problem:} Find a general solution to the differential equation $y'' - 3y' + 2y = e^x \sin x$.

\textbf{Solution:}
\textbf{Step 1: Find the homogeneous solution, $y_h(x)$.}
The associated homogeneous equation is $y'' - 3y' + 2y = 0$.
The characteristic equation is:
\[ r^2 - 3r + 2 = 0 \]
Factoring the quadratic:
\[ (r - 1)(r - 2) = 0 \implies r_1 = 1, \quad r_2 = 2 \]
Thus, the homogeneous solution is:
\[ y_h(x) = c_1 e^x + c_2 e^{2x} \]

\textbf{Step 2: Determine the form of the particular solution, $y_p(x)$.}
The forcing function is $f(x) = e^x \sin x$. Here, $\alpha = 1$ and $\beta = 1$.
The complex root associated with this forcing function is $\alpha + i\beta = 1 + i$.
We check if $1 + i$ is a root of our characteristic equation. The roots are $1$ and $2$. Since $1 + i$ is not a root, the multiplicity is $s = 0$.
Our guess for the particular solution is:
\[ y_p(x) = e^x (A \cos x + B \sin x) \]

\textbf{Step 3: Differentiate $y_p(x)$ and substitute into the ODE.}
Using the product rule:
\begin{align*}
y_p' &= e^x (A \cos x + B \sin x) + e^x (-A \sin x + B \cos x) \\
     &= e^x \left[ (A + B) \cos x + (B - A) \sin x \right]
\end{align*}
Differentiating again:
\begin{align*}
y_p'' &= e^x \left[ (A + B) \cos x + (B - A) \sin x \right] + e^x \left[ -(A + B) \sin x + (B - A) \cos x \right] \\
      &= e^x \left[ (A + B + B - A) \cos x + (B - A - A - B) \sin x \right] \\
      &= e^x \left[ 2B \cos x - 2A \sin x \right]
\end{align*}
Substitute $y_p, y_p', y_p''$ into $y'' - 3y' + 2y = e^x \sin x$:
\[ e^x (2B \cos x - 2A \sin x) - 3e^x \left[ (A + B) \cos x + (B - A) \sin x \right] + 2e^x (A \cos x + B \sin x) = e^x \sin x \]
Divide the entire equation by $e^x$ (since $e^x \neq 0$):
\[ (2B \cos x - 2A \sin x) - (3A + 3B) \cos x - (3B - 3A) \sin x + 2A \cos x + 2B \sin x = \sin x \]
Group the $\cos x$ and $\sin x$ terms:
\[ \cos x \left[ 2B - 3A - 3B + 2A \right] + \sin x \left[ -2A - 3B + 3A + 2B \right] = \sin x \]
Simplify the coefficients:
\[ \cos x \left[ -A - B \right] + \sin x \left[ A - B \right] = 0 \cos x + 1 \sin x \]

\textbf{Step 4: Solve for A and B.}
Equate the coefficients on both sides:
\begin{align*}
-A - B &= 0 \implies B = -A \\
A - B &= 1
\end{align*}
Substitute $B = -A$ into the second equation:
\[ A - (-A) = 1 \implies 2A = 1 \implies A = \frac{1}{2} \]
Since $B = -A$, we have $B = -\frac{1}{2}$.
Thus, the particular solution is:
\[ y_p(x) = e^x \left( \frac{1}{2} \cos x - \frac{1}{2} \sin x \right) \]

\textbf{Step 5: State the general solution.}
\[ y(x) = y_h(x) + y_p(x) \]
\textbf{Final Answer:} $y(x) = c_1 e^x + c_2 e^{2x} + \frac{1}{2} e^x (\cos x - \sin x)$

\subsection{Problem 28: Initial Value Problem with Multiple Forcing Terms}
\textbf{Problem:} Find the solution to the initial value problem: $y'' + y' - 12y = e^t + e^{2t} - 1$; $y(0) = 1, y'(0) = 3$.

\textbf{Solution:}
\textbf{Step 1: Find the homogeneous solution.}
Characteristic equation: $r^2 + r - 12 = 0 \implies (r + 4)(r - 3) = 0$.
Roots are $r = -4, 3$.
\[ y_h(t) = c_1 e^{-4t} + c_2 e^{3t} \]

\textbf{Step 2: Find the particular solution using Superposition.}
The forcing function is $f(t) = e^t + e^{2t} - 1$. We split this into three parts: $f_1(t) = e^t$, $f_2(t) = e^{2t}$, $f_3(t) = -1$. None of the roots of these forcing functions ($1, 2, 0$) overlap with the homogeneous roots ($-4, 3$).
\begin{itemize}
    \item For $f_1(t) = e^t$, guess $y_{p1} = A e^t$.
    \[ A e^t + A e^t - 12 A e^t = e^t \implies -10A = 1 \implies A = -\frac{1}{10} \]
    \item For $f_2(t) = e^{2t}$, guess $y_{p2} = B e^{2t}$.
    \[ 4B e^{2t} + 2B e^{2t} - 12B e^{2t} = e^{2t} \implies -6B = 1 \implies B = -\frac{1}{6} \]
    \item For $f_3(t) = -1$, guess $y_{p3} = C$.
    \[ 0 + 0 - 12C = -1 \implies C = \frac{1}{12} \]
\end{itemize}
The total particular solution is $y_p(t) = -\frac{1}{10}e^t - \frac{1}{6}e^{2t} + \frac{1}{12}$.
The general solution is:
\[ y(t) = c_1 e^{-4t} + c_2 e^{3t} - \frac{1}{10}e^t - \frac{1}{6}e^{2t} + \frac{1}{12} \]

\textbf{Step 3: Apply Initial Conditions.}
First, find $y'(t)$:
\[ y'(t) = -4c_1 e^{-4t} + 3c_2 e^{3t} - \frac{1}{10}e^t - \frac{2}{6}e^{2t} \]
Apply $y(0) = 1$:
\[ c_1 + c_2 - \frac{1}{10} - \frac{1}{6} + \frac{1}{12} = 1 \]
Find a common denominator (60) for the fractions: $-\frac{6}{60} - \frac{10}{60} + \frac{5}{60} = -\frac{11}{60}$.
\[ c_1 + c_2 - \frac{11}{60} = 1 \implies c_1 + c_2 = \frac{71}{60} \]
Apply $y'(0) = 3$:
\[ -4c_1 + 3c_2 - \frac{1}{10} - \frac{1}{3} = 3 \]
Common denominator (30): $-\frac{3}{30} - \frac{10}{30} = -\frac{13}{30}$.
\[ -4c_1 + 3c_2 - \frac{13}{30} = 3 \implies -4c_1 + 3c_2 = \frac{103}{30} = \frac{206}{60} \]
We have a system of equations:
\begin{align*}
1) \quad c_1 + c_2 &= \frac{71}{60} \\
2) \quad -4c_1 + 3c_2 &= \frac{206}{60}
\end{align*}
Multiply equation (1) by 4: $4c_1 + 4c_2 = \frac{284}{60}$.
Add to equation (2):
\[ 7c_2 = \frac{284 + 206}{60} = \frac{490}{60} \implies c_2 = \frac{70}{60} = \frac{7}{6} \]
Substitute $c_2$ back into (1):
\[ c_1 + \frac{70}{60} = \frac{71}{60} \implies c_1 = \frac{1}{60} \]

\textbf{Final Answer:} $y(t) = \frac{1}{60} e^{-4t} + \frac{7}{6} e^{3t} - \frac{1}{10}e^t - \frac{1}{6}e^{2t} + \frac{1}{12}$

\subsection{Problem 31: Determining the Form of $y_p$ (Resonance and Tricks)}
\textbf{Problem:} Determine the form of a particular solution for $y'' + y = \sin t + t \cos t + 10^t$. Do not solve.

\textbf{Solution:}
\textbf{Step 1: Analyze the homogeneous equation.}
$r^2 + 1 = 0 \implies r = \pm i$.
The homogeneous roots are $0 \pm 1i$.

\textbf{Step 2: Analyze the forcing function pieces.}
The forcing function is $f(t) = \sin t + t \cos t + 10^t$.
\begin{itemize}
    \item \textbf{Piece 1:} $f_1(t) = \sin t + t \cos t$. This is associated with $\alpha = 0, \beta = 1$, so the root is $\pm i$. Because $\pm i$ is a single root of the characteristic equation, we have a collision (resonance). Thus, $s = 1$. The polynomial attached to the trig functions is of degree 1 ($t$).
    The base guess would be $(At + B)\cos t + (Ct + D)\sin t$.
    Multiplying by $t^1$ due to resonance gives: $t(At + B)\cos t + t(Ct + D)\sin t$.
    \item \textbf{Piece 2:} $f_2(t) = 10^t$. This looks unusual, but we can rewrite it using base $e$: $10^t = (e^{\ln 10})^t = e^{(\ln 10)t}$. This is a standard exponential function where $\alpha = \ln 10$. Since $\ln 10$ is not a root of the characteristic equation, $s = 0$.
    The guess is simply $E e^{(\ln 10)t}$ or $E 10^t$.
\end{itemize}

\textbf{Step 3: Combine the guesses.}
\textbf{Final Answer:} $y_p(t) = t(At + B)\cos t + t(Ct + D)\sin t + E 10^t$

\subsection{Problem 43: Real-World Modeling (Mass-Spring System)}
\textbf{Problem:} A mass-spring system is driven by a sinusoidal external force $g(t) = 5 \sin t$. The mass equals 1, the spring constant equals 3, and the damping coefficient equals 4. If the mass is initially located at $y(0) = 1/2$ and at rest ($y'(0) = 0$), find its equation of motion.

\textbf{Solution:}
\textbf{Step 1: Set up the differential equation.}
The standard model is $m y'' + b y' + k y = g(t)$.
Substituting the given values: $y'' + 4y' + 3y = 5 \sin t$.

\textbf{Step 2: Find the homogeneous solution (Transient State).}
$r^2 + 4r + 3 = 0 \implies (r + 1)(r + 3) = 0 \implies r = -1, -3$.
$y_h(t) = c_1 e^{-t} + c_2 e^{-3t}$.

\textbf{Step 3: Find the particular solution (Steady-State).}
Guess $y_p(t) = A \cos t + B \sin t$.
$y_p'(t) = -A \sin t + B \cos t$.
$y_p''(t) = -A \cos t - B \sin t$.
Substitute into the DE:
\[ (-A \cos t - B \sin t) + 4(-A \sin t + B \cos t) + 3(A \cos t + B \sin t) = 5 \sin t \]
Group terms:
\[ \cos t (-A + 4B + 3A) + \sin t (-B - 4A + 3B) = 5 \sin t \]
\[ \cos t (2A + 4B) + \sin t (-4A + 2B) = 5 \sin t \]
Equate coefficients:
\begin{align*}
2A + 4B &= 0 \implies A = -2B \\
-4A + 2B &= 5
\end{align*}
Substitute $A$: $-4(-2B) + 2B = 5 \implies 10B = 5 \implies B = 1/2$.
Then $A = -2(1/2) = -1$.
$y_p(t) = -\cos t + \frac{1}{2} \sin t$.

\textbf{Step 4: Apply Initial Conditions.}
$y(t) = c_1 e^{-t} + c_2 e^{-3t} - \cos t + \frac{1}{2} \sin t$.
$y(0) = c_1 + c_2 - 1 = \frac{1}{2} \implies c_1 + c_2 = \frac{3}{2}$.
$y'(t) = -c_1 e^{-t} - 3c_2 e^{-3t} + \sin t + \frac{1}{2} \cos t$.
$y'(0) = -c_1 - 3c_2 + \frac{1}{2} = 0 \implies c_1 + 3c_2 = \frac{1}{2}$.
Subtract the second equation from the first:
\[ (c_1 + c_2) - (c_1 + 3c_2) = \frac{3}{2} - \frac{1}{2} \implies -2c_2 = 1 \implies c_2 = -\frac{1}{2} \]
Substitute $c_2$ into the first equation: $c_1 - \frac{1}{2} = \frac{3}{2} \implies c_1 = 2$.

\textbf{Final Answer:} $y(t) = 2e^{-t} - \frac{1}{2}e^{-3t} - \cos t + \frac{1}{2}\sin t$

\subsection{Problem 48: Reverse Engineering a Differential Equation}
\textbf{Problem:} All that is known concerning a mysterious second-order constant-coefficient differential equation $y'' + py' + qy = g(t)$ is that $t^2 + 1 + e^t \cos t$, $t^2 + 1 + e^t \sin t$, and $t^2 + 1 + e^t \cos t + e^t \sin t$ are solutions.
(a) Determine two linearly independent solutions to the corresponding homogeneous equation.
(b) Find a suitable choice of $p, q$, and $g(t)$.

\textbf{Solution:}
\textbf{Part (a):}
Let the three given solutions be $y_1, y_2, y_3$. By the theory of linear ODEs, the difference between any two solutions of a nonhomogeneous equation is a solution to the corresponding homogeneous equation.
\[ y_{h1} = y_2 - y_1 = (t^2 + 1 + e^t \sin t) - (t^2 + 1 + e^t \cos t) = e^t \sin t - e^t \cos t \]
\[ y_{h2} = y_3 - y_1 = (t^2 + 1 + e^t \cos t + e^t \sin t) - (t^2 + 1 + e^t \cos t) = e^t \sin t \]
\[ y_{h3} = y_3 - y_2 = (t^2 + 1 + e^t \cos t + e^t \sin t) - (t^2 + 1 + e^t \sin t) = e^t \cos t \]
From this, we can clearly see that the fundamental set of solutions to the homogeneous equation is $\{e^t \cos t, e^t \sin t\}$.

\textbf{Part (b):}
Since the homogeneous solutions are $e^t \cos t$ and $e^t \sin t$, the roots of the characteristic equation must be complex conjugates: $r = 1 \pm i$.
We can reconstruct the characteristic polynomial:
\[ (r - (1 + i))(r - (1 - i)) = (r - 1)^2 - i^2 = r^2 - 2r + 1 - (-1) = r^2 - 2r + 2 = 0 \]
Thus, the homogeneous equation is $y'' - 2y' + 2y = 0$. This means $p = -2$ and $q = 2$.
To find $g(t)$, we note that $y_1 = (t^2 + 1) + e^t \cos t$. Since $e^t \cos t$ is $y_h$, the remaining part $y_p = t^2 + 1$ must be a particular solution to the nonhomogeneous equation.
We plug $y_p = t^2 + 1$ into the left side of our reconstructed DE to find $g(t)$:
$y_p' = 2t$
$y_p'' = 2$
\[ g(t) = y_p'' - 2y_p' + 2y_p = 2 - 2(2t) + 2(t^2 + 1) = 2 - 4t + 2t^2 + 2 = 2t^2 - 4t + 4 \]
\textbf{Final Answer:} (a) $e^t \cos t$ and $e^t \sin t$. (b) $p = -2, q = 2, g(t) = 2t^2 - 4t + 4$.

\newpage
\section{Part 3: In-Depth Analysis of Problems and Techniques}

\subsection{A) Problem Types and General Approach}
\begin{itemize}
    \item \textbf{Superposition Application (Problem 1):} These problems test your understanding of linearity. \textit{Approach:} Do not solve the DE. Simply multiply the given particular solutions by the same constants that multiply the forcing functions in the new equation, and add them together.
    \item \textbf{Applicability Checks (Problem 11):} These test the limitations of the method. \textit{Approach:} Scan the forcing function for terms that are not polynomials, exponentials, or sines/cosines. Terms like $1/t, \tan t, \ln t$ immediately disqualify the method of undetermined coefficients.
    \item \textbf{Full General Solution (Problem 19):} The standard bread-and-butter problem. \textit{Approach:} Always find $y_h$ first. Then formulate a guess for $y_p$, differentiate twice, substitute, and solve the resulting system of equations.
    \item \textbf{Initial Value Problems (Problem 28):} \textit{Approach:} Find the full general solution $y(t) = y_h + y_p$ \textbf{before} applying any initial conditions. A common fatal error is applying initial conditions to $y_h$ alone.
    \item \textbf{Form Determination / Resonance (Problem 31):} \textit{Approach:} Find the roots of the characteristic equation. Compare them to the "roots" implied by the forcing function. If they match, multiply your standard guess by $t^s$, where $s$ is the multiplicity of the root in the homogeneous equation.
    \item \textbf{Reverse Engineering (Problem 48):} \textit{Approach:} Use the property that $y_a - y_b = y_h$ to isolate the homogeneous solutions. Rebuild the characteristic equation from the roots to find the left side of the DE. Plug the remaining nonhomogeneous part into the left side to reveal the forcing function.
\end{itemize}

\subsection{B) Key Algebraic and Calculus Manipulations}
\begin{itemize}
    \item \textbf{The $10^t$ Trick (Problem 31):} Recognizing that $a^t = e^{t \ln a}$ is crucial. The method of undetermined coefficients is built around base $e$. Converting bases allows the method to proceed normally.
    \item \textbf{Grouping by Linear Independence (Problem 19 \& 43):} After substituting $y_p$ into the DE, you are left with a massive string of terms. The critical algebraic step is factoring out the linearly independent functions (e.g., grouping all $\cos x$ terms together and all $\sin x$ terms together). This transforms one differential equation into a system of algebraic equations.
    \item \textbf{Product Rule Efficiency (Problem 19):} When differentiating $e^{\alpha x}(A \cos \beta x + B \sin \beta x)$, factoring out the $e^{\alpha x}$ at each step prevents the expression from exploding into unmanageable fragments.
\end{itemize}

\newpage
\section{Part 4: "Cheatsheet" and Tips for Success}

\subsection{Summary of Important Formulas}
\begin{tcolorbox}
\textbf{The Master Guess Table for $y_p(t)$:}
\begin{itemize}
    \item If $f(t) = P_n(t)$ (Polynomial of degree $n$): \\ Guess $y_p = t^s (A_n t^n + \dots + A_1 t + A_0)$
    \item If $f(t) = P_n(t) e^{\alpha t}$: \\ Guess $y_p = t^s (A_n t^n + \dots + A_0) e^{\alpha t}$
    \item If $f(t) = P_n(t) e^{\alpha t} \cos(\beta t)$ or $P_n(t) e^{\alpha t} \sin(\beta t)$: \\ Guess $y_p = t^s \left[ (A_n t^n + \dots + A_0) e^{\alpha t} \cos(\beta t) + (B_n t^n + \dots + B_0) e^{\alpha t} \sin(\beta t) \right]$
\end{itemize}
\textit{* $s$ is the number of times $\alpha + i\beta$ is a root of the characteristic equation.}
\end{tcolorbox}

\subsection{Cheats, Tricks, and Shortcuts}
\begin{itemize}
    \item \textbf{The "Zero" Check:} Before doing any heavy algebra, plug $t=0$ into your $y_p, y_p', y_p''$ expressions if possible to do a quick sanity check on your coefficients.
    \item \textbf{Divide and Conquer:} If $f(t)$ has 4 different terms, do not try to solve a massive $8 \times 8$ system of equations. Break it into 4 separate $y_p$ problems, solve them individually, and add the results. It is much less prone to arithmetic errors.
    \item \textbf{Missing Terms:} If $f(t) = \sin t$, your guess MUST include both $\sin t$ and $\cos t$. Never assume a coefficient is zero just because the term isn't in the forcing function.
\end{itemize}

\subsection{Common Pitfalls and Traps}
\begin{itemize}
    \item \textbf{The Resonance Trap:} Forgetting to multiply by $t^s$ when the forcing function matches the homogeneous solution. If you forget this, your system of equations will yield a contradiction (like $0 = 1$).
    \item \textbf{The IVP Trap:} Applying $y(0)$ and $y'(0)$ to $y_h(t)$ instead of the full $y(t) = y_h(t) + y_p(t)$. The initial conditions apply to the \textit{entire} physical system, not just the transient part.
    \item \textbf{The "Close Enough" Root Trap:} If $y_h = c_1 \cos(2t) + c_2 \sin(2t)$ and $f(t) = \cos(3t)$, there is NO resonance. The roots must match exactly ($\pm 2i \neq \pm 3i$). Do not multiply by $t$.
\end{itemize}

\newpage
\section{Part 5: Conceptual Synthesis and The "Big Picture"}

\subsection{A) Thematic Connections}
The central theme of this topic is \textbf{Linearity and Decomposition}. In mathematics, when a system is linear, the whole is exactly equal to the sum of its parts. We saw this in Calculus I (the derivative of a sum is the sum of derivatives) and in Linear Algebra (vector transformations). Here, it manifests as the Superposition Principle. We are taking a complex, messy reality (a system subjected to multiple different forces) and decomposing it into isolated, simple scenarios that we can solve independently.

\subsection{B) Forward and Backward Links}
\textbf{Backward Links:} This topic is the immediate logical extension of solving homogeneous linear ODEs. The characteristic equation ($r^2 + br + c = 0$) you learned previously is now the diagnostic tool used to determine the "s" multiplier (resonance) for the nonhomogeneous case.

\textbf{Forward Links:} The Method of Undetermined Coefficients is powerful but fragile—it breaks the moment $f(t) = \sec t$ or $f(t) = \ln t$. This limitation directly motivates the next major topic: \textbf{Variation of Parameters}. Furthermore, the concept of breaking a forcing function into simpler parts lays the groundwork for \textbf{Fourier Series} and \textbf{Laplace Transforms}, where we will represent \textit{any} periodic function as an infinite sum of sines and cosines, allowing us to solve ODEs for virtually any forcing function.

\newpage
\section{Part 6: Real-World Application and Modeling}

\subsection{A) Concrete Scenarios in Quantitative Finance and Economics}
While traditionally applied to physics (mass-spring systems, RLC circuits), second-order linear ODEs with forcing terms are increasingly vital in advanced financial modeling and macroeconomics.

\begin{enumerate}
    \item \textbf{Macroeconomic Business Cycle Modeling:} Economists model the fluctuations of Gross Domestic Product (GDP) around a long-term growth trend using second-order ODEs. The "restoring force" represents market corrections, and the "damping" represents monetary policy smoothing. The forcing function $f(t)$ represents external deterministic shocks, such as a cyclical government stimulus package or seasonal agricultural outputs.
    \item \textbf{Commodity Spread Mean Reversion:} In quantitative trading, the price spread between two correlated commodities (like Brent and WTI crude oil) is often modeled as a mean-reverting process with momentum (requiring a second derivative). The forcing function $f(t)$ can represent known, deterministic seasonal demand cycles (e.g., higher heating oil demand in winter), modeled as a sinusoidal wave.
\end{enumerate}

\subsection{B) Model Problem Setup: Commodity Spread Pricing}
Let $S(t)$ represent the price spread between two assets at time $t$. We hypothesize that the spread exhibits momentum (mass) and mean-reversion (spring), subject to a seasonal cyclical demand shock.

\textbf{Variables:}
\begin{itemize}
    \item $S(t)$: The price spread at time $t$.
    \item $\gamma$: The rate of liquidity dissipation (analogous to damping/friction).
    \item $\kappa$: The strength of mean-reversion (analogous to spring constant).
    \item $F_0$: The amplitude of the seasonal demand shock.
    \item $\omega$: The frequency of the seasonal cycle (e.g., $\omega = 2\pi/12$ for monthly data in a yearly cycle).
\end{itemize}

\textbf{The Mathematical Model:}
The dynamics of the spread can be modeled by the nonhomogeneous second-order differential equation:
\[ S''(t) + \gamma S'(t) + \kappa S(t) = F_0 \cos(\omega t) \]
To find the expected trajectory of the spread $S(t)$, a quantitative analyst would use the Method of Undetermined Coefficients to find the steady-state seasonal oscillation $S_p(t)$, and add it to the transient mean-reverting path $S_h(t)$.

\newpage
\section{Part 7: Common Variations and Untested Concepts}

Based on standard differential equations curricula, there are two important concepts related to this chapter that were not explicitly tested in the provided homework set.

\subsection{Untested Concept 1: Higher Multiplicity of Roots ($s=2$)}
The homework featured problems where the forcing function's root matched the homogeneous root once ($s=1$). However, if the characteristic equation has a \textit{repeated} root that also matches the forcing function, you must multiply by $t^2$.

\textbf{Example:} Find the form of $y_p$ for $y'' - 2y' + y = 5e^t$.
\textbf{Explanation:} 
The homogeneous equation is $r^2 - 2r + 1 = 0 \implies (r-1)^2 = 0$. The root is $r=1$ with multiplicity 2.
The forcing function is $5e^t$, which corresponds to $\alpha = 1$.
Because $\alpha=1$ is a \textit{double} root of the characteristic equation, $s=2$.
If we guessed $y_p = A e^t$, it would vanish (it's in $y_h$).
If we guessed $y_p = A t e^t$, it would also vanish (it's also in $y_h$).
Therefore, the correct form is:
\[ y_p(t) = A t^2 e^t \]

\subsection{Untested Concept 2: When Undetermined Coefficients Fails (Preview of Variation of Parameters)}
The homework asked you to identify \textit{when} the method fails (Problem 11 with $1/t$), but it didn't ask you to solve such a problem. When the forcing function is a quotient of trig functions (like $\sec t$ or $\tan t$) or involves logarithms, we must use a different method called Variation of Parameters.

\textbf{Example:} $y'' + y = \sec t$
\textbf{Explanation:} 
The forcing function is $\sec t$. If we try to guess $y_p = A \sec t$, the derivative is $A \sec t \tan t$. The second derivative introduces $\sec^3 t$ and $\sec t \tan^2 t$. The derivatives never loop back to a finite set of terms. Undetermined coefficients is impossible here. In the next section, you will learn to set $y_p = u_1(t)y_1(t) + u_2(t)y_2(t)$ and solve for the functions $u_1$ and $u_2$ using integration.

\newpage
\section{Part 8: Advanced Diagnostic Testing: "Find the Flaw"}
\textit{Instructions: You are the grader. Each of the following five problems contains a complete solution with a single, subtle, yet critical flaw. Find the error, explain why it is an error in one sentence, and provide the correct final solution.}

\subsection{Problem 1}
\textbf{Solve:} $y'' - 4y = e^{2t}$
\textbf{Flawed Solution:}
1. $r^2 - 4 = 0 \implies r = \pm 2$. So $y_h = c_1 e^{2t} + c_2 e^{-2t}$.
2. Guess $y_p = A e^{2t}$.
3. $y_p' = 2A e^{2t}$, $y_p'' = 4A e^{2t}$.
4. Substitute: $4A e^{2t} - 4(A e^{2t}) = e^{2t} \implies 0 = e^{2t}$.
5. Since $0 \neq e^{2t}$, there is no particular solution.
\textbf{Your Task:} Find the flaw and correct it.

\subsection{Problem 2}
\textbf{Solve:} $y'' + 9y = \sin(3t)$
\textbf{Flawed Solution:}
1. $r^2 + 9 = 0 \implies r = \pm 3i$. $y_h = c_1 \cos(3t) + c_2 \sin(3t)$.
2. Forcing is $\sin(3t)$, root is $3i$. This matches $y_h$, so $s=1$.
3. Guess $y_p = A t \sin(3t)$.
4. $y_p' = A \sin(3t) + 3At \cos(3t)$.
5. $y_p'' = 3A \cos(3t) + 3A \cos(3t) - 9At \sin(3t) = 6A \cos(3t) - 9At \sin(3t)$.
6. Substitute: $6A \cos(3t) - 9At \sin(3t) + 9(A t \sin(3t)) = \sin(3t)$.
7. $6A \cos(3t) = \sin(3t)$. This is impossible, so the method fails.
\textbf{Your Task:} Find the flaw and correct it.

\subsection{Problem 3}
\textbf{Solve:} $y'' - y' - 2y = t^2$
\textbf{Flawed Solution:}
1. $r^2 - r - 2 = 0 \implies (r-2)(r+1) = 0 \implies y_h = c_1 e^{2t} + c_2 e^{-t}$.
2. Forcing is $t^2$. Guess $y_p = A t^2$.
3. $y_p' = 2At$, $y_p'' = 2A$.
4. Substitute: $2A - (2At) - 2(A t^2) = t^2$.
5. Equate coefficients: $-2A = 1 \implies A = -1/2$.
6. Final answer: $y = c_1 e^{2t} + c_2 e^{-t} - \frac{1}{2}t^2$.
\textbf{Your Task:} Find the flaw and correct it.

\subsection{Problem 4}
\textbf{Solve:} $y'' + y = e^t \cos t$
\textbf{Flawed Solution:}
1. $r^2 + 1 = 0 \implies r = \pm i$. $y_h = c_1 \cos t + c_2 \sin t$.
2. Forcing is $e^t \cos t$. Because $\cos t$ is in the forcing function and in $y_h$, we have resonance.
3. We must multiply by $t$. Guess $y_p = t e^t (A \cos t + B \sin t)$.
4. (Proceeds with massive product rule derivatives...)
\textbf{Your Task:} Find the flaw and correct it.

\subsection{Problem 5}
\textbf{Solve:} $y'' - 2y' + y = 4e^t$
\textbf{Flawed Solution:}
1. $r^2 - 2r + 1 = 0 \implies (r-1)^2 = 0 \implies r = 1$ (double root).
2. $y_h = c_1 e^t + c_2 t e^t$.
3. Forcing is $4e^t$. Root is $1$. This matches the homogeneous root, so we multiply by $t$.
4. Guess $y_p = A t e^t$.
5. $y_p' = A e^t + A t e^t$, $y_p'' = 2A e^t + A t e^t$.
6. Substitute: $(2A e^t + A t e^t) - 2(A e^t + A t e^t) + A t e^t = 4e^t$.
7. $2A e^t - 2A e^t = 4e^t \implies 0 = 4e^t$. No solution.
\textbf{Your Task:} Find the flaw and correct it.

\end{document}